% GENERATED FILE. Edit canonical lessons, then run npm run generate:books.
% canonical-source-hash: fnv1a64:c54a0a1999680ee5
% canonical-lessons: LA-C429-carrus, LA-C429-cisium, LA-C429-essedum, LA-C429-carpentum, LA-C429-carruca

\chapter{Wagon, Light gig, War chariot, Covered carriage, Travelling coach}
\label{ch:la-wagon-light-gig-war-chariot-covered-carriage-travelling-coach}
\hlchaptermodality{429}

\begin{chapteropening}
\textbf{By the end of this chapter:} \emph{I can name Roman vehicles: a wagon, a light gig, a two-wheeled chariot, a covered carriage and a travelling coach.}

Complete the last lesson of chapter 429: I can name Roman vehicles: a wagon, a light gig, a two-wheeled chariot, a covered carriage and a travelling coach.
\end{chapteropening}

\section[carrus, carrī]{carrus, carrī — a wagon, a baggage cart}

\label{lesson:LA-C429-carrus}



\begin{quote}
Before the new one: say the Latin for to travel abroad, then the Latin for to go forward.
\end{quote}

\subsection*{You'll want to know: carrus, carrī}
\textbf{carrus, carrī} — \textquotedblleft{}a wagon, a baggage cart\textquotedblright{}.

\begin{quote}
Carrus frūmentum ex agrīs ad oppidum vehit.
\end{quote}

\begin{quote}
\emph{The wagon carries grain from the fields to the town.}
\end{quote}

Second declension, masculine: \emph{carrus, carrī m.} (also \emph{carrum, carrī n.}). A heavy cart for goods. English \emph{car}, \emph{carry} and \emph{career} come from it.

\begin{cousinweb}
\emph{carrus} was borrowed from Gaulish \emph{karros} 'wagon', from PIE \emph{\'{k}ers-} 'to run' (EtymDB); Lewis \& Short give no origin. English \emph{car}, \emph{carry} and \emph{career} go back to it.
\end{cousinweb}

\begin{grammarlens}[title={What you've built}]
One more word for telling, explaining and describing at B1.
\end{grammarlens}

\subsection*{Your turn}
\begin{itemize}
\raggedright
  \item \emph{Say it:} \emph{carrus, carrī}
  \item \emph{Say it:} \emph{carrus, carrī}, once more
  \item \emph{Say it:} a sentence of your own with \emph{carrus, carrī}
  \item \emph{Recall:} say the Latin for to travel abroad, then the Latin for to go forward, then say \emph{carrus, carrī} again
\end{itemize}

\begin{tcolorbox}[breakable,colback=teal!4,colframe=teal!35!black,title={Before you move on}]
What does carrus, carrī mean? (\textquotedblleft{}A wagon, a baggage cart\textquotedblright{}.) Say it once more.
\end{tcolorbox}
\section[cisium, cisiī]{cisium, cisiī — a light gig, a fast two-wheeled cart}

\label{lesson:LA-C429-cisium}



\begin{quote}
Before the new one: say the Latin for to go forward, then the Latin for a wagon.
\end{quote}

\subsection*{You'll want to know: cisium, cisiī}
\textbf{cisium, cisiī} — \textquotedblleft{}a light gig, a fast two-wheeled cart\textquotedblright{}.

\begin{quote}
Cisium celere est; itaque mercātor cisiō ad urbem it.
\end{quote}

\begin{quote}
\emph{A gig is fast; so the merchant goes to the city by gig.}
\end{quote}

Second declension, neuter: \emph{cisium, cisiī n.} A light two-wheeled carriage for one or two people, for speed rather than baggage.

\begin{cousinweb}
The origin of \emph{cisium} is not settled. Lewis \& Short give the sense, a light two-wheeled vehicle, but no source, and EtymDB does not list it.
\end{cousinweb}

\begin{grammarlens}[title={What you've built}]
One more word for telling, explaining and describing at B1.
\end{grammarlens}

\subsection*{Your turn}
\begin{itemize}
\raggedright
  \item \emph{Say it:} \emph{cisium, cisiī}
  \item \emph{Say it:} \emph{cisium, cisiī}, once more
  \item \emph{Say it:} a sentence of your own with \emph{cisium, cisiī}
  \item \emph{Recall:} say the Latin for to go forward, then the Latin for a wagon, then say \emph{cisium, cisiī} again
\end{itemize}

\begin{tcolorbox}[breakable,colback=teal!4,colframe=teal!35!black,title={Before you move on}]
What does cisium, cisiī mean? (\textquotedblleft{}A light gig, a fast two-wheeled cart\textquotedblright{}.) Say it once more.
\end{tcolorbox}
\section[essedum, essedī]{essedum, essedī — a two-wheeled chariot; a light travelling carriage}

\label{lesson:LA-C429-essedum}



\begin{quote}
Before the new one: say the Latin for a wagon, then the Latin for a light gig.
\end{quote}

\subsection*{You'll want to know: essedum, essedī}
\textbf{essedum, essedī} — \textquotedblleft{}a two-wheeled chariot; a light travelling carriage\textquotedblright{}.

\begin{quote}
Dīves essedō celerī ad vīllam vehitur.
\end{quote}

\begin{quote}
\emph{The rich man rides to his country house in a fast two-wheeled carriage.}
\end{quote}

Second declension, neuter: \emph{essedum, essedī n.} First a war chariot of northern peoples, later a smart carriage in the city.

\begin{cousinweb}
\emph{essedum} names the war chariot of the Gauls and Britons (Lewis \& Short, under \emph{essedārius}), later a travelling carriage. It is usually thought to be a Gaulish word, but neither Lewis \& Short nor EtymDB states an origin, so the line is hedged.
\end{cousinweb}

\begin{grammarlens}[title={What you've built}]
One more word for telling, explaining and describing at B1.
\end{grammarlens}

\subsection*{Your turn}
\begin{itemize}
\raggedright
  \item \emph{Say it:} \emph{essedum, essedī}
  \item \emph{Say it:} \emph{essedum, essedī}, once more
  \item \emph{Say it:} a sentence of your own with \emph{essedum, essedī}
  \item \emph{Recall:} say the Latin for a wagon, then the Latin for a light gig, then say \emph{essedum, essedī} again
\end{itemize}

\begin{tcolorbox}[breakable,colback=teal!4,colframe=teal!35!black,title={Before you move on}]
What does essedum, essedī mean? (\textquotedblleft{}A two-wheeled chariot; a light travelling carriage\textquotedblright{}.) Say it once more.
\end{tcolorbox}
\section[carpentum, carpentī]{carpentum, carpentī — a covered carriage; a two-wheeled coach}

\label{lesson:LA-C429-carpentum}



\begin{quote}
Before the new one: say the Latin for a light gig, then the Latin for a two-wheeled chariot.
\end{quote}

\subsection*{You'll want to know: carpentum, carpentī}
\textbf{carpentum, carpentī} — \textquotedblleft{}a covered carriage; a two-wheeled coach\textquotedblright{}.

\begin{quote}
Māter et fīliae carpentō ad vīllam vehuntur.
\end{quote}

\begin{quote}
\emph{The mother and her daughters ride to the country house in a covered carriage.}
\end{quote}

Second declension, neuter: \emph{carpentum, carpentī n.} Married women of rank rode in it. English \emph{carpenter} first meant a maker of carriages.

\begin{cousinweb}
\emph{carpentum} was borrowed from Gaulish \emph{carbantos} 'two-wheeled carriage' (EtymDB); Lewis \& Short give no origin. English \emph{carpenter} comes, through French, from Latin \emph{carpentārius} 'carriage-maker', built on it.
\end{cousinweb}

\begin{grammarlens}[title={What you've built}]
One more word for telling, explaining and describing at B1.
\end{grammarlens}

\subsection*{Your turn}
\begin{itemize}
\raggedright
  \item \emph{Say it:} \emph{carpentum, carpentī}
  \item \emph{Say it:} \emph{carpentum, carpentī}, once more
  \item \emph{Say it:} a sentence of your own with \emph{carpentum, carpentī}
  \item \emph{Recall:} say the Latin for a light gig, then the Latin for a two-wheeled chariot, then say \emph{carpentum, carpentī} again
\end{itemize}

\begin{tcolorbox}[breakable,colback=teal!4,colframe=teal!35!black,title={Before you move on}]
What does carpentum, carpentī mean? (\textquotedblleft{}A covered carriage; a two-wheeled coach\textquotedblright{}.) Say it once more.
\end{tcolorbox}
\section[carrūca, carrūcae]{carrūca, carrūcae — a travelling coach, a four-wheeled carriage}

\label{lesson:LA-C429-carruca}



\begin{quote}
Before the new one: say the Latin for a two-wheeled chariot, then the Latin for a covered carriage.
\end{quote}

\subsection*{You'll want to know: carrūca, carrūcae}
\textbf{carrūca, carrūcae} — \textquotedblleft{}a travelling coach, a four-wheeled carriage\textquotedblright{}.

\begin{quote}
Dīves in carrūcā sedet et servī post eum ambulant.
\end{quote}

\begin{quote}
\emph{The rich man sits in his coach and the slaves walk behind him.}
\end{quote}

First declension, feminine: \emph{carrūca, carrūcae f.} A grand four-wheeled coach. In French it later became \emph{charrue}, a plough.

\begin{cousinweb}
\emph{carrūca} is built on \emph{carrus} 'wagon' (EtymDB), which was borrowed from Gaulish \emph{karros}. Lewis \& Short gloss it as a four-wheeled travelling and state coach.
\end{cousinweb}

\begin{grammarlens}[title={What you've built}]
One more word for telling, explaining and describing at B1.
\end{grammarlens}

\subsection*{Your turn}
\begin{itemize}
\raggedright
  \item \emph{Say it:} \emph{carrūca, carrūcae}
  \item \emph{Say it:} \emph{carrūca, carrūcae}, once more
  \item \emph{Say it:} a sentence of your own with \emph{carrūca, carrūcae}
  \item \emph{Recall:} say the Latin for a two-wheeled chariot, then the Latin for a covered carriage, then say \emph{carrūca, carrūcae} again
\end{itemize}

\begin{tcolorbox}[breakable,colback=teal!4,colframe=teal!35!black,title={Before you move on}]
What does carrūca, carrūcae mean? (\textquotedblleft{}A travelling coach, a four-wheeled carriage\textquotedblright{}.) Say it once more.
\end{tcolorbox}
